\section{Tensor product}

Let $R$ be a commutative ring, and let $M$ and $N$ be $R$-modules. Roughly speaking, the tensor product $M \otimes_R N$ linearizes bilinear maps out of $M \times N$.

\begin{definition}
    Let $R$ be a commutative ring, and let $M$, $N$, and $K$ be $R$-modules. A map
    \[
        f\colon M \times N \to K
    \]
    is called \emph{bilinear} if for all $r\in R$, $m,m_1,m_2\in M$, and $n,n_1,n_2\in N$, one has
    \begin{itemize}
        \item $f(m_1+m_2,n) = f(m_1,n) + f(m_2,n)$;
        \item $f(m,n_1+n_2) = f(m,n_1) + f(m,n_2)$;
        \item $f(rm,n) = r f(m,n)$;
        \item $f(m,rn) = r f(m,n)$.
    \end{itemize}
\end{definition}

\begin{example}
    If $f\colon M\times N \to K$ is bilinear, then
    \[
        f(x,0)=0=f(0,y)
    \]
    for all $x\in M$ and $y\in N$. Moreover,
    \[
        f(-x,y)=-f(x,y), \qquad f(x,-y)=-f(x,y)
    \]
    for all $x\in M$ and $y\in N$.
\end{example}

\begin{definition}
    The \emph{tensor product} of $M$ and $N$ is a pair $(M \otimes_R N, t_{M,N})$, where $M \otimes_R N$ is an $R$-module and
    \[
        t_{M,N}\colon M\times N \to M \otimes_R N
    \]
    is a bilinear map, usually written $t_{M,N}(m,n)=m\otimes n$, such that the following universal property holds:

    whenever $X$ is an $R$-module and $f\colon M\times N \to X$ is bilinear, there exists a unique $R$-module homomorphism
    \[
        \tilde{f}\colon M \otimes_R N \to X
    \]
    such that the following diagram commutes:
    \begin{center}
        \begin{tikzcd}
            M \times N \arrow[r, "f"] \arrow[d, "{t_{M,N}}"'] & X \\
            M \otimes_R N \arrow[ru, "\tilde{f}"', dashed]    &
        \end{tikzcd}
    \end{center}
\end{definition}

\begin{theorem}
    The tensor product of $M$ and $N$ exists and is unique up to isomorphism.
\end{theorem}

\begin{proof}
    We first prove uniqueness.

    Suppose $(T_1,t_1)$ and $(T_2,t_2)$ are two tensor products of $M$ and $N$. By the universal property, there exist unique $R$-module homomorphisms
    \[
        \varphi\colon T_1 \to T_2,
        \qquad
        \psi\colon T_2 \to T_1
    \]
    such that
    \[
        t_2 = \varphi \circ t_1,
        \qquad
        t_1 = \psi \circ t_2.
    \]
    Hence
    \[
        \varphi \circ \psi \circ t_2 = \varphi \circ t_1 = t_2,
        \qquad
        \psi \circ \varphi \circ t_1 = \psi \circ t_2 = t_1.
    \]
    By uniqueness in the universal property, it follows that
    \[
        \varphi \circ \psi = \operatorname{id}_{T_2},
        \qquad
        \psi \circ \varphi = \operatorname{id}_{T_1}.
    \]
    Therefore $\varphi$ and $\psi$ are inverse isomorphisms, so $T_1 \cong T_2$.

    Now we prove existence.

    Let $F$ be the free $R$-module with basis $M \times N$, say
    \[
        F = \bigoplus_{(m,n)\in M\times N} R e_{m,n}.
    \]
    Let $B$ be the submodule of $F$ generated by the following elements:
    \begin{itemize}
        \item $e_{m_1+m_2,n}-e_{m_1,n}-e_{m_2,n}$;
        \item $e_{m,n_1+n_2}-e_{m,n_1}-e_{m,n_2}$;
        \item $e_{rm,n}-r e_{m,n}$;
        \item $e_{m,rn}-r e_{m,n}$,
    \end{itemize}
    where $m,m_1,m_2\in M$, $n,n_1,n_2\in N$, and $r\in R$.

    Define
    \[
        M\otimes_R N = F/B,
        \qquad
        t_{M,N}(m,n)=e_{m,n}+B.
    \]
    Then $t_{M,N}\colon M\times N \to M\otimes_R N$ is a well-defined bilinear map.

    It remains to verify the universal property.

    Let $X$ be an $R$-module, and let $f\colon M\times N \to X$ be bilinear. Since $F$ is free on the set $M\times N$, there exists a unique $R$-module homomorphism
    \[
        \hat{f}\colon F \to X
    \]
    such that $\hat{f}(e_{m,n}) = f(m,n)$ for all $(m,n)\in M\times N$.

    We claim that $B \subseteq \ker \hat{f}$. Indeed, since $f$ is bilinear,
    \begin{align*}
        \hat{f}(e_{m_1+m_2,n}-e_{m_1,n}-e_{m_2,n}) & = f(m_1+m_2,n)-f(m_1,n)-f(m_2,n) = 0, \\
        \hat{f}(e_{m,n_1+n_2}-e_{m,n_1}-e_{m,n_2}) & = f(m,n_1+n_2)-f(m,n_1)-f(m,n_2) = 0, \\
        \hat{f}(e_{rm,n}-r e_{m,n})                & = f(rm,n)-r f(m,n) = 0,               \\
        \hat{f}(e_{m,rn}-r e_{m,n})                & = f(m,rn)-r f(m,n) = 0.
    \end{align*}
    Since $B$ is generated by these elements and $\ker \hat{f}$ is a submodule, we obtain $B \subseteq \ker \hat{f}$.

    By the universal property of quotient modules, $\hat{f}$ induces a unique $R$-module homomorphism
    \[
        \tilde{f}\colon M\otimes_R N = F/B \to X
    \]
    such that $\tilde{f} \circ \pi = \hat{f}$, where $\pi\colon F \to F/B$ is the canonical projection. In particular,
    \[
        \tilde{f}(m\otimes n)
        =
        \tilde{f}(e_{m,n}+B)
        =
        \hat{f}(e_{m,n})
        =
        f(m,n),
    \]
    so $\tilde{f} \circ t_{M,N} = f$.

    For uniqueness, suppose $g\colon M\otimes_R N \to X$ is another $R$-module homomorphism such that $g \circ t_{M,N} = f$. Then
    \[
        g(m\otimes n) = f(m,n) = \tilde{f}(m\otimes n)
    \]
    for all $m\in M$ and $n\in N$. Since the elementary tensors generate $M\otimes_R N$ as an $R$-module, it follows that $g=\tilde{f}$.

    Therefore $(M\otimes_R N, t_{M,N})$ is a tensor product of $M$ and $N$.
\end{proof}

\begin{definition}
    The set of elements of the form $m \otimes n$ for $m \in M$ and $n \in N$ is called the set of elementary tensors. This set generates $M \otimes_R N$ as an $R$-module.
\end{definition}

\begin{example}
    Fix $m \in M$. The map $N \to M \otimes_R N, n \mapsto m \otimes n$ is a morphism of $R$-modules. Its image $m \otimes_R N = \{m \otimes n \mid n \in N\}$ is a submodule of $M \otimes_R N$ which is isomorphic to a quotient $N/\operatorname{Ann}_R(m)N$.
\end{example}

\begin{proposition} We have the following isomorphisms:
    \begin{enumerate}
        \item $(M_1 \otimes_R M_2) \otimes_R M_3 \cong M_1 \otimes_R (M_2 \otimes_R M_3)$ where $(m_1 \otimes m_2) \otimes m_3 \mapsto m_1 \otimes (m_2 \otimes m_3)$.
        \item $M_1 \otimes_R M_2 \cong M_2 \otimes_R M_1$ where $m_1 \otimes m_2 \mapsto m_2 \otimes m_1$.
        \item $M_1 \otimes_R (M_2 \oplus M_3) \cong (M_1 \otimes_R M_2) \oplus (M_1 \otimes_R M_3)$ where $m_1 \otimes (m_2 \oplus m_3) \mapsto (m_1 \otimes m_2) + (m_1 \otimes m_3)$. And more generally,
              \[
                  M_1 \otimes_R (\bigoplus_{i=1}^n M_i) \cong \bigoplus_{i=1}^n (M_1 \otimes_R M_i)
              \]
              where $m_1 \otimes (m_2 \oplus \cdots \oplus m_n) \mapsto (m_1 \otimes m_2) + \cdots + (m_1 \otimes m_n)$.
        \item $R \otimes_R M \cong M$ where $r \otimes m \mapsto rm$. And more generally,$R^{(I)} \otimes_R M \cong M^{(I)}$ where $(r_i)_{i\in I} \otimes m \mapsto (r_i m)_{i\in I}$.
    \end{enumerate}
\end{proposition}

\begin{proof}
    We prove (1):$(M_1 \otimes_R M_2) \otimes_R M_3 \cong M_1 \otimes_R (M_2 \otimes_R M_3)$.

    \textbf{Step 1.} We construct $\varphi: (M_1 \otimes_R M_2) \otimes_R M_3 \to M_1 \otimes_R (M_2 \otimes_R M_3)$.

    Fix $m_3 \in M_3$. For each $m_3$, define $f_{m_3}: M_1 \times M_2 \to M_1 \otimes_R (M_2 \otimes_R M_3)$ by $f_{m_3}(m_1, m_2) = m_1 \otimes (m_2 \otimes m_3)$. This is bilinear in $(m_1, m_2)$, so by the universal property of $M_1 \otimes_R M_2$, there exists a unique $R$-module homomorphism $\tilde{f}_{m_3}: M_1 \otimes_R M_2 \to M_1 \otimes_R (M_2 \otimes_R M_3)$ such that $\tilde{f}_{m_3}(m_1 \otimes m_2) = m_1 \otimes (m_2 \otimes m_3)$:
    \begin{center}
        \begin{tikzcd}
            M_1 \times M_2 \arrow[r, "f_{m_3}"] \arrow[d, "t_{M_1,M_2}"'] & M_1 \otimes_R (M_2 \otimes_R M_3) \\
            M_1 \otimes_R M_2 \arrow[ru, "\tilde{f}_{m_3}"', dashed] &
        \end{tikzcd}
    \end{center}

    Now define $g: (M_1 \otimes_R M_2) \times M_3 \to M_1 \otimes_R (M_2 \otimes_R M_3)$ by $g(\alpha, m_3) = \tilde{f}_{m_3}(\alpha)$. We verify that $g$ is bilinear. Linearity in the first argument follows because each $\tilde{f}_{m_3}$ is an $R$-module homomorphism. For the second argument, on elementary tensors:
    \begin{align*}
        g(m_1 \otimes m_2, m_3 + m_3') & = m_1 \otimes (m_2 \otimes (m_3+m_3'))                           \\
                                       & = m_1 \otimes (m_2 \otimes m_3 + m_2 \otimes m_3')               \\
                                       & = m_1 \otimes (m_2 \otimes m_3) + m_1 \otimes (m_2 \otimes m_3') \\
                                       & = g(m_1 \otimes m_2, m_3) + g(m_1 \otimes m_2, m_3'),            \\
        g(m_1 \otimes m_2, rm_3)       & = m_1 \otimes (m_2 \otimes rm_3)                                 \\
                                       & = m_1 \otimes (r(m_2 \otimes m_3))                               \\
                                       & = r(m_1 \otimes (m_2 \otimes m_3))                               \\
                                       & = r\, g(m_1 \otimes m_2, m_3).
    \end{align*}
    Since elementary tensors generate $M_1 \otimes_R M_2$, bilinearity holds on all of $(M_1 \otimes_R M_2) \times M_3$.

    By the universal property of $(M_1 \otimes_R M_2) \otimes_R M_3$, there exists a unique $R$-module homomorphism
    \[
        \varphi: (M_1 \otimes_R M_2) \otimes_R M_3 \to M_1 \otimes_R (M_2 \otimes_R M_3)
    \]
    such that $\varphi((m_1 \otimes m_2) \otimes m_3) = m_1 \otimes (m_2 \otimes m_3)$:
    \begin{center}
        \begin{tikzcd}
            (M_1 \otimes_R M_2) \times M_3 \arrow[r, "g"] \arrow[d, "t"'] & M_1 \otimes_R (M_2 \otimes_R M_3) \\
            (M_1 \otimes_R M_2) \otimes_R M_3 \arrow[ru, "\varphi"', dashed] &
        \end{tikzcd}
    \end{center}

    \textbf{Step 2.} By an entirely symmetric argument (fixing $m_1$ first, then extending), we obtain an $R$-module homomorphism
    \[
        \psi: M_1 \otimes_R (M_2 \otimes_R M_3) \to (M_1 \otimes_R M_2) \otimes_R M_3
    \]
    such that $\psi(m_1 \otimes (m_2 \otimes m_3)) = (m_1 \otimes m_2) \otimes m_3$.

    \textbf{Step 3.} We show $\varphi$ and $\psi$ are inverses. On elementary tensors:
    \[
        (\psi \circ \varphi)((m_1 \otimes m_2) \otimes m_3) = \psi(m_1 \otimes (m_2 \otimes m_3)) = (m_1 \otimes m_2) \otimes m_3,
    \]
    \[
        (\varphi \circ \psi)(m_1 \otimes (m_2 \otimes m_3)) = \varphi((m_1 \otimes m_2) \otimes m_3) = m_1 \otimes (m_2 \otimes m_3).
    \]
    Since elementary tensors generate both modules,$\psi \circ \varphi = \operatorname{id}$ and $\varphi \circ \psi = \operatorname{id}$. The overall situation is summarized by the following commutative diagram:
    \begin{center}
        \begin{tikzcd}
            & M_1 \times M_2 \times M_3 \arrow[ld] \arrow[rd] & \\
            (M_1 \otimes_R M_2) \otimes_R M_3 \arrow[rr, "\varphi", shift left=1.5] & & M_1 \otimes_R (M_2 \otimes_R M_3) \arrow[ll, "\psi", shift left=1.5]
        \end{tikzcd}
    \end{center}
    where the diagonal arrows send $(m_1,m_2,m_3)$ to $(m_1\otimes m_2)\otimes m_3$ and $m_1\otimes(m_2\otimes m_3)$ respectively. Hence $\varphi$ is an isomorphism with $(m_1\otimes m_2)\otimes m_3 \mapsto m_1\otimes(m_2\otimes m_3)$.
\end{proof}

\begin{theorem}
    Let $M, N, P$ be $R$-modules. We have a canonical isomorphism:
    \[
        \phi_{M,N,P}: \operatorname{Hom}_R(M \otimes_R N, P) \to \operatorname{Hom}_R(M, \operatorname{Hom}_R(N, P))
    \]
    such that $\phi_{M,N,P}(f)(m)(n) = f(m \otimes n)$.
    Furthermore, if $R$ is commutative, then $\phi_{M,N,P}$ is an $R$-module isomorphism. If $R$ is not commutative, then $\phi_{M,N,P}$ is only a $\mathbb{Z}$-module isomorphism.
\end{theorem}

\begin{proof}
    Let $f \in \operatorname{Hom}_R(M \otimes_R N, P)$. For each $m \in M$, define $\phi_{M,N,P}(f)(m): N \to P$ by
    \[
        \phi_{M,N,P}(f)(m)(n) = f(m \otimes n).
    \]
    We verify that $\phi_{M,N,P}(f)(m) \in \operatorname{Hom}_R(N, P)$: for $n_1, n_2 \in N$ and $r \in R$,
    \begin{align*}
        \phi_{M,N,P}(f)(m)(n_1 + n_2) & = f(m \otimes (n_1+n_2)) = f(m \otimes n_1 + m \otimes n_2) \\
                                      & = f(m\otimes n_1) + f(m\otimes n_2),                        \\
        \phi_{M,N,P}(f)(m)(rn)        & = f(m \otimes rn) = f(r(m\otimes n))                        \\
                                      & = rf(m\otimes n) = r\,\phi_{M,N,P}(f)(m)(n).
    \end{align*}
    Next, we verify that $\phi_{M,N,P}(f) \in \operatorname{Hom}_R(M, \operatorname{Hom}_R(N,P))$: for $m_1, m_2 \in M$,$r \in R$, and any $n\in N$,
    \begin{align*}
        \phi_{M,N,P}(f)(m_1+m_2)(n) & = f((m_1+m_2)\otimes n)             \\
                                    & = f(m_1\otimes n) + f(m_2\otimes n) \\
        \phi_{M,N,P}(f)(rm)(n)      & = f(rm\otimes n)                    \\
                                    & = f(r(m\otimes n))                  \\
                                    & = rf(m\otimes n)                    \\
                                    & = (r\,\phi_{M,N,P}(f)(m))(n).
    \end{align*}
    So
    \[
        \phi_{M,N,P}(f)(m_1+m_2) = \phi_{M,N,P}(f)(m_1)+\phi_{M,N,P}(f)(m_2)
    \]
    and $\phi_{M,N,P}(f)(rm) = r\,\phi_{M,N,P}(f)(m)$, confirming $R$-linearity.

    Let $g \in \operatorname{Hom}_R(M, \operatorname{Hom}_R(N, P))$. Define $\hat{g}: M \times N \to P$ by $\hat{g}(m,n) = g(m)(n)$. We check that $\hat{g}$ is $R$-bilinear:
    \begin{align*}
        \hat{g}(m_1+m_2,n) & = g(m_1+m_2)(n) = g(m_1)(n)+g(m_2)(n) = \hat{g}(m_1,n)+\hat{g}(m_2,n), \\
        \hat{g}(m,n_1+n_2) & = g(m)(n_1+n_2) = g(m)(n_1)+g(m)(n_2) = \hat{g}(m,n_1)+\hat{g}(m,n_2), \\
        \hat{g}(rm,n)      & = g(rm)(n) = (rg(m))(n) = rg(m)(n) = r\,\hat{g}(m,n),                  \\
        \hat{g}(m,rn)      & = g(m)(rn) = rg(m)(n) = r\,\hat{g}(m,n).
    \end{align*}
    By the universal property of $M \otimes_R N$, there exists a unique $R$-module homomorphism $\psi_{M,N,P}(g): M\otimes_R N \to P$ such that $\psi_{M,N,P}(g)(m\otimes n) = g(m)(n)$:
    \begin{center}
        \begin{tikzcd}
            M \times N \arrow[r, "\hat{g}"] \arrow[d, "t_{M,N}"'] & P \\
            M \otimes_R N \arrow[ru, "{\psi_{M,N,P}(g)}"', dashed] &
        \end{tikzcd}
    \end{center}

    For any $f \in \operatorname{Hom}_R(M\otimes_R N,P)$:
    \[
        \psi_{M,N,P}(\phi_{M,N,P}(f))(m\otimes n) = \phi_{M,N,P}(f)(m)(n) = f(m\otimes n).
    \]
    Since elementary tensors generate $M\otimes_R N$, we have $\psi_{M,N,P}(\phi_{M,N,P}(f)) = f$.

    For any $g \in \operatorname{Hom}_R(M, \operatorname{Hom}_R(N,P))$, for all $m\in M$ and $n\in N$:
    \[
        \phi_{M,N,P}(\psi_{M,N,P}(g))(m)(n) = \psi_{M,N,P}(g)(m\otimes n) = g(m)(n).
    \]
    So $\phi_{M,N,P}(\psi_{M,N,P}(g))(m) = g(m)$ for all $m$, hence $\phi_{M,N,P}(\psi_{M,N,P}(g)) = g$.

    Therefore $\phi_{M,N,P}$ is a bijection of sets (in fact, of abelian groups).

    It remains to show that $\phi_{M,N,P}$ respects the appropriate module structure.

    Both $\operatorname{Hom}_R(M\otimes_R N, P)$ and $\operatorname{Hom}_R(M, \operatorname{Hom}_R(N,P))$ are abelian groups under pointwise addition. We verify $\phi_{M,N,P}$ is a group homomorphism: for $f_1, f_2 \in \operatorname{Hom}_R(M\otimes_R N, P)$,
    \begin{align*}
        \phi_{M,N,P}(f_1+f_2)(m)(n) & = (f_1+f_2)(m\otimes n)                           \\
                                    & = f_1(m\otimes n)+f_2(m\otimes n)                 \\
                                    & = \phi_{M,N,P}(f_1)(m)(n)+\phi_{M,N,P}(f_2)(m)(n)
    \end{align*}
    so $\phi_{M,N,P}(f_1+f_2)=\phi_{M,N,P}(f_1)+\phi_{M,N,P}(f_2)$. Hence $\phi_{M,N,P}$ is a $\mathbb{Z}$-module isomorphism.

    If $R$ is commutative,$\operatorname{Hom}_R(M\otimes_R N, P)$ carries an $R$-module structure via $(r\cdot f)(x) = rf(x)$, and similarly for $\operatorname{Hom}_R(M, \operatorname{Hom}_R(N,P))$. Then
    \begin{align*}
        \phi_{M,N,P}(rf)(m)(n) & = (rf)(m\otimes n)           \\
                               & = r f(m\otimes n)            \\
                               & = r\,\phi_{M,N,P}(f)(m)(n)   \\
                               & = (r\,\phi_{M,N,P}(f))(m)(n)
    \end{align*}
    so $\phi_{M,N,P}(rf) = r\,\phi_{M,N,P}(f)$, and $\phi_{M,N,P}$ is an $R$-module isomorphism.

    If $R$ is not commutative,$\operatorname{Hom}_R(M\otimes_R N, P)$ does not naturally carry an $R$-module structure (scalar multiplication by $r$ need not commute with $R$-linearity), so $\phi_{M,N,P}$ is only a $\mathbb{Z}$-module isomorphism.
\end{proof}

\begin{definition}
    Let $f: M\to M'$ and $g: N\to N'$ be $R$-module homomorphisms. The tensor product of $f$ and $g$ is the $R$-module homomorphism:
    \[
        f \otimes g: M \otimes_R N \to M' \otimes_R N'
    \]
    defined by $(f \otimes g)(m \otimes n) = f(m) \otimes g(n)$.
\end{definition}

\begin{definition}
    Let $M_i, i \in I$ be a family of $R$-modules. Use the following notation:
    \[
        L^n(M_1, \ldots, M_n; N) = \{ f: M_1 \times \cdots \times M_n \to N \mid f \text{ is $n$-linear } \}
    \]
    In particular,
    \[
        L^2(M_1, M_2; N) \cong L(M_1 \otimes M_2, N) \cong \operatorname{Hom}_R(M_1 \otimes M_2, N)
    \]
    or equivalently,
    \[
        L^2(M_1 \otimes M_2, N) \cong L(M_1, L(M_2, N))
    \]
\end{definition}

\begin{proposition}
    Let $M,N$ be two $R$-modules where $M$ is finitely generated and free. Denote
    \[
        M^{\vee} = \operatorname{Hom}_R(M, R)
    \]
    Then we have a canonical isomorphism:
    \[
        M^{\vee} \otimes N \cong \operatorname{Hom}_R(M, N)
    \]
\end{proposition}

\begin{proof}
    Since $M$ is a finitely generated free $R$-module, we have $M \cong R^m$ for some $m \geq 1$. We construct an explicit isomorphism
    \[
        \Phi: M^{\vee} \otimes_R N \to \operatorname{Hom}_R(M, N).
    \]
    Define a bilinear map $\beta: M^{\vee} \times N \to \operatorname{Hom}_R(M,N)$ by
    \[
        \beta(f, n)(m) = f(m)\, n, \qquad f \in M^{\vee},\; n \in N,\; m \in M.
    \]
    We check $R$-bilinearity. For $f_1, f_2 \in M^{\vee}$,$n, n_1, n_2 \in N$,$m \in M$, and $r \in R$:
    \begin{align*}
        \beta(f_1+f_2, n)(m) & = (f_1+f_2)(m)\,n = f_1(m)\,n + f_2(m)\,n     \\
                             & = \beta(f_1,n)(m)+\beta(f_2,n)(m),            \\
        \beta(f, n_1+n_2)(m) & = f(m)(n_1+n_2) = f(m)\,n_1+f(m)\,n_2         \\
                             & = \beta(f,n_1)(m)+\beta(f,n_2)(m),            \\
        \beta(rf, n)(m)      & = (rf)(m)\,n = r\,f(m)\,n = r\,\beta(f,n)(m), \\
        \beta(f, rn)(m)      & = f(m)\,(rn) = r\,f(m)\,n = r\,\beta(f,n)(m).
    \end{align*}
    By the universal property of the tensor product, there exists a unique $R$-module homomorphism
    \[
        \Phi: M^{\vee} \otimes_R N \to \operatorname{Hom}_R(M, N), \qquad \Phi(f \otimes n)(m) = f(m)\,n.
    \]

    Now let $\{e_1, \ldots, e_m\}$ be a basis of $M \cong R^m$, and let $\{e_1^*, \ldots, e_m^*\}$ be the dual basis; that is,
    \[
        e_i^* \in M^{\vee},
        \qquad
        e_i^*(e_j) = \delta_{ij}.
    \]
    Define
    \[
        \Psi: \operatorname{Hom}_R(M, N) \to M^{\vee} \otimes_R N, \qquad \Psi(g) = \sum_{i=1}^{m} e_i^* \otimes g(e_i).
    \]
    We verify $\Psi$ is an $R$-module homomorphism: for $g_1, g_2 \in \operatorname{Hom}_R(M,N)$ and $r \in R$,
    \begin{align*}
        \Psi(g_1+g_2) & = \sum_{i=1}^m e_i^* \otimes (g_1+g_2)(e_i) = \sum_{i=1}^m e_i^* \otimes (g_1(e_i)+g_2(e_i))       \\
                      & = \sum_{i=1}^m e_i^* \otimes g_1(e_i) + \sum_{i=1}^m e_i^* \otimes g_2(e_i) = \Psi(g_1)+\Psi(g_2), \\
        \Psi(rg)      & = \sum_{i=1}^m e_i^* \otimes (rg)(e_i) = \sum_{i=1}^m e_i^* \otimes r\,g(e_i)                      \\
                      & = r \sum_{i=1}^m e_i^* \otimes g(e_i) = r\,\Psi(g).
    \end{align*}

    We show $\Phi$ and $\Psi$ are mutual inverses. For any $g \in \operatorname{Hom}_R(M,N)$ and $m = \sum_j r_j e_j \in M$:
    \begin{align*}
        \Phi(\Psi(g))(m) & = \Phi\!\Bigl(\sum_{i=1}^m e_i^* \otimes g(e_i)\Bigr)\!\Bigl(\sum_j r_j e_j\Bigr)   \\
                         & = \sum_{i=1}^m e_i^*\!\Bigl(\sum_j r_j e_j\Bigr) g(e_i) = \sum_{i=1}^m r_i\, g(e_i) \\
                         & = g\!\Bigl(\sum_{i=1}^m r_i e_i\Bigr) = g(m).
    \end{align*}
    So $\Phi \circ \Psi = \operatorname{id}$.

    For any elementary tensor $f \otimes n \in M^{\vee} \otimes_R N$:
    \begin{align*}
        \Psi(\Phi(f \otimes n)) & = \sum_{i=1}^m e_i^* \otimes \Phi(f \otimes n)(e_i) = \sum_{i=1}^m e_i^* \otimes f(e_i)\,n                   \\
                                & = \sum_{i=1}^m f(e_i)\, (e_i^* \otimes n) = \Bigl(\sum_{i=1}^m f(e_i)\, e_i^*\Bigr) \otimes n = f \otimes n,
    \end{align*}
    where the last equality uses $f = \sum_{i=1}^m f(e_i)\, e_i^*$(expansion of $f$ in the dual basis). Since elementary tensors generate $M^{\vee} \otimes_R N$ and $\Psi \circ \Phi$ is $R$-linear, we have $\Psi \circ \Phi = \operatorname{id}$.

    Therefore $\Phi$ is an $R$-module isomorphism, giving $M^{\vee} \otimes_R N \cong \operatorname{Hom}_R(M, N)$.
\end{proof}

\begin{lemma}
    Let $P$ be a finitely generated projective $R$-module. Then there exists an
    $R$-module $Q$ and an integer $t\geq 0$ such that
    \[
        P\oplus Q\cong R^t.
    \]
\end{lemma}

\begin{proof}
    Since $P$ is finitely generated, there exist elements
    \[
        p_1,\dots,p_t\in P
    \]
    such that
    \[
        P=Rp_1+\cdots+Rp_t.
    \]
    Therefore we have a surjective $R$-module homomorphism
    \[
        \pi:R^t\longrightarrow P
    \]
    defined by
    \[
        \pi(a_1,\dots,a_t)=a_1p_1+\cdots+a_tp_t.
    \]
    Equivalently, if $e_1,\dots,e_t$ denotes the standard basis of $R^t$, then
    \[
        \pi(e_i)=p_i
        \qquad
        (1\leq i\leq t).
    \]
    Since $p_1,\dots,p_t$ generate $P$, the map $\pi$ is surjective.

    Now we use that $P$ is projective. Since $\pi:R^t\to P$ is an epimorphism and
    $P$ is projective, the map $\pi$ splits. That is, there exists an $R$-module
    homomorphism
    \[
        s:P\longrightarrow R^t
    \]
    such that
    \[
        \pi\circ s=\operatorname{id}_P.
    \]
    In particular, $s$ is injective, because if $s(p)=0$, then
    \[
        p=(\pi\circ s)(p)=\pi(0)=0.
    \]

    We claim that
    \[
        R^t=\operatorname{Im}(s)\oplus \ker\pi.
    \]
    First, let $x\in R^t$. Then
    \[
        x=s(\pi(x))+\bigl(x-s(\pi(x))\bigr).
    \]
    Clearly,
    \[
        s(\pi(x))\in \operatorname{Im}(s).
    \]
    Moreover,
    \[
        \pi\bigl(x-s(\pi(x))\bigr)
        =
        \pi(x)-(\pi\circ s)(\pi(x))
        =
        \pi(x)-\pi(x)
        =
        0.
    \]
    Hence
    \[
        x-s(\pi(x))\in \ker\pi.
    \]
    Therefore
    \[
        R^t=\operatorname{Im}(s)+\ker\pi.
    \]

    It remains to show that the sum is direct. Suppose
    \[
        y\in \operatorname{Im}(s)\cap \ker\pi.
    \]
    Then $y=s(p)$ for some $p\in P$, and also $\pi(y)=0$. Hence
    \[
        0=\pi(y)=\pi(s(p))=(\pi\circ s)(p)=p.
    \]
    Therefore
    \[
        y=s(p)=s(0)=0.
    \]
    Thus
    \[
        \operatorname{Im}(s)\cap \ker\pi=0.
    \]
    Hence
    \[
        R^t=\operatorname{Im}(s)\oplus \ker\pi.
    \]

    Since $s:P\to \operatorname{Im}(s)$ is an isomorphism, we obtain
    \[
        R^t\cong P\oplus \ker\pi.
    \]
    Taking
    \[
        Q:=\ker\pi,
    \]
    we get
    \[
        P\oplus Q\cong R^t.
    \]
\end{proof}

\begin{remark}
    This does not mean that $P$ itself is free. It only says that $P$ is a direct
    summand of a finite free module. Thus a finitely generated projective module is
    ``part of'' a finite free module:
    \[
        R^t\cong P\oplus Q.
    \]
\end{remark}

\begin{corollary}
    Let $R$ be a commutative ring with $1$, let $P$ be a finitely generated projective
    $R$-module, and let $N$ be any $R$-module. Denote
    \[
        P^{\vee}:=\operatorname{Hom}_R(P,R).
    \]
    Then the canonical $R$-module homomorphism
    \[
        \tau_{P,N}:P^{\vee}\otimes_R N\longrightarrow \operatorname{Hom}_R(P,N)
    \]
    defined by
    \[
        \tau_{P,N}(\varphi\otimes n)(p)=\varphi(p)n
        \qquad
        (\varphi\in P^\vee,\ n\in N,\ p\in P)
    \]
    is an isomorphism.
\end{corollary}

\begin{proof}
    Since $P$ is finitely generated and projective, there exist elements
    \[
        p_1,\dots,p_t\in P
    \]
    and homomorphisms
    \[
        \varphi_1,\dots,\varphi_t\in P^\vee
    \]
    such that
    \[
        p=\sum_{i=1}^t \varphi_i(p)p_i
        \qquad
        \text{for every }p\in P.
    \]
    Indeed, since $P$ is finitely generated projective, there exists an $R$-module $Q$
    and an integer $t\geq 0$ such that
    \[
        P\oplus Q\cong R^t.
    \]
    Let
    \[
        \iota:P\longrightarrow R^t
        \qquad\text{and}\qquad
        \rho:R^t\longrightarrow P
    \]
    be the corresponding inclusion and projection, so that
    \[
        \rho\circ \iota=\operatorname{id}_P.
    \]
    Let $e_1,\dots,e_t$ be the standard basis of $R^t$, and let
    $e_1^\vee,\dots,e_t^\vee$ be the dual basis of $(R^t)^\vee$. Define
    \[
        p_i:=\rho(e_i)\in P,
        \qquad
        \varphi_i:=e_i^\vee\circ \iota\in P^\vee.
    \]
    Then, for every $p\in P$,
    \[
        p
        =
        \rho(\iota(p))
        =
        \rho\left(\sum_{i=1}^t e_i^\vee(\iota(p))e_i\right)
        =
        \sum_{i=1}^t \varphi_i(p)p_i.
    \]

    Now define a map
    \[
        \sigma_{P,N}:\operatorname{Hom}_R(P,N)\longrightarrow P^\vee\otimes_R N
    \]
    by
    \[
        \sigma_{P,N}(f)
        =
        \sum_{i=1}^t \varphi_i\otimes f(p_i).
    \]
    We show that $\sigma_{P,N}$ is the inverse of $\tau_{P,N}$.

    First, let $f\in \operatorname{Hom}_R(P,N)$. For every $p\in P$, we have
    \[
        \bigl(\tau_{P,N}\circ \sigma_{P,N}\bigr)(f)(p)
        =
        \tau_{P,N}\left(\sum_{i=1}^t \varphi_i\otimes f(p_i)\right)(p)
        =
        \sum_{i=1}^t \varphi_i(p)f(p_i).
    \]
    Since $f$ is $R$-linear,
    \[
        \sum_{i=1}^t \varphi_i(p)f(p_i)
        =
        f\left(\sum_{i=1}^t \varphi_i(p)p_i\right)
        =
        f(p).
    \]
    Hence
    \[
        \tau_{P,N}\circ \sigma_{P,N}
        =
        \operatorname{id}_{\operatorname{Hom}_R(P,N)}.
    \]

    Conversely, let $\psi\otimes n\in P^\vee\otimes_R N$ be an elementary tensor.
    Then
    \[
        \bigl(\sigma_{P,N}\circ \tau_{P,N}\bigr)(\psi\otimes n)
        =
        \sigma_{P,N}\bigl(\tau_{P,N}(\psi\otimes n)\bigr)
        =
        \sum_{i=1}^t
        \varphi_i\otimes \tau_{P,N}(\psi\otimes n)(p_i).
    \]
    By definition of $\tau_{P,N}$,
    \[
        \tau_{P,N}(\psi\otimes n)(p_i)
        =
        \psi(p_i)n.
    \]
    Therefore
    \[
        \bigl(\sigma_{P,N}\circ \tau_{P,N}\bigr)(\psi\otimes n)
        =
        \sum_{i=1}^t \varphi_i\otimes \psi(p_i)n
        =
        \left(\sum_{i=1}^t \psi(p_i)\varphi_i\right)\otimes n.
    \]
    But for every $p\in P$,
    \[
        \left(\sum_{i=1}^t \psi(p_i)\varphi_i\right)(p)
        =
        \sum_{i=1}^t \psi(p_i)\varphi_i(p)
        =
        \sum_{i=1}^t \varphi_i(p)\psi(p_i)
        =
        \psi\left(\sum_{i=1}^t \varphi_i(p)p_i\right)
        =
        \psi(p).
    \]
    Hence
    \[
        \sum_{i=1}^t \psi(p_i)\varphi_i=\psi.
    \]
    Thus
    \[
        \bigl(\sigma_{P,N}\circ \tau_{P,N}\bigr)(\psi\otimes n)
        =
        \psi\otimes n.
    \]
    Since elementary tensors generate $P^\vee\otimes_R N$, it follows that
    \[
        \sigma_{P,N}\circ \tau_{P,N}
        =
        \operatorname{id}_{P^\vee\otimes_R N}.
    \]

    Therefore $\tau_{P,N}$ is an isomorphism, with inverse $\sigma_{P,N}$.
\end{proof}

\begin{proposition}
    Assume that $R$ is a subring of $S$, and $M$ is an $R$-module. Then the bilinear map:
    \[
        S \times (S \otimes_R M) \to S \otimes_R M
    \]
    defined by
    \[
        (s, s' \otimes m) \mapsto s s' \otimes m
    \]
    defines a structure of $S$-module on the $R$-module $S \otimes_R M$.

    More generally, let $f: R \to T$ be a ring homomorphism. And view $T$ as an $R$-module via $R \times T \to T, (\lambda, \mu) \mapsto f(\lambda) \mu$. Then the bilinear map:
    \[
        T \times (T \otimes_R M) \to T \otimes_R M
    \]
    defined by
    \[
        (t, t' \otimes m) \mapsto t t' \otimes m
    \]
    defines a structure of $T$-module on the $R$-module $T \otimes_R M$.
\end{proposition}

\begin{proof}
    It suffices to prove the more general statement in \textup{(2)}, since \textup{(1)} is the
    special case where $T=S$ and $f:R\hookrightarrow S$ is the inclusion map.

    For each fixed $t\in T$, define a map
    \[
        \beta_t:T\times M \longrightarrow T\otimes_R M,\qquad (t',m)\longmapsto tt'\otimes m.
    \]
    This map is $R$-bilinear. Hence, by the universal property of the tensor product,
    there exists a unique $R$-linear map
    \[
        \mu_t:T\otimes_R M\longrightarrow T\otimes_R M
    \]
    such that
    \[
        \mu_t(t'\otimes m)=tt'\otimes m
        \qquad\text{for all }t'\in T,\ m\in M.
    \]

    Now define
    \[
        t\cdot x:=\mu_t(x),\qquad t\in T,\ x\in T\otimes_R M.
    \]
    We claim that this gives a left $T$-module structure on $T\otimes_R M$.

    Since the elementary tensors $t'\otimes m$ generate $T\otimes_R M$ as an abelian
    group, it is enough to verify the module axioms on elementary tensors.

    Let $t_1,t_2,t'\in T$ and $m\in M$. Then
    \[
        (t_1+t_2)\cdot (t'\otimes m)
        =(t_1+t_2)t'\otimes m
        =t_1t'\otimes m+t_2t'\otimes m
        =t_1\cdot (t'\otimes m)+t_2\cdot (t'\otimes m).
    \]
    Also,
    \[
        (t_1t_2)\cdot (t'\otimes m)
        =t_1t_2t'\otimes m
        =t_1\cdot (t_2t'\otimes m)
        =t_1\cdot \bigl(t_2\cdot (t'\otimes m)\bigr).
    \]
    Moreover,
    \[
        1_T\cdot (t'\otimes m)=1_Tt'\otimes m=t'\otimes m.
    \]
    Finally, for any $x,y\in T\otimes_R M$,
    \[
        t\cdot (x+y)=\mu_t(x+y)=\mu_t(x)+\mu_t(y)=t\cdot x+t\cdot y,
    \]
    because $\mu_t$ is $R$-linear, hence additive.

    Therefore the operation
    \[
        T\times (T\otimes_R M)\longrightarrow T\otimes_R M,\qquad
        (t,\,t'\otimes m)\longmapsto tt'\otimes m
    \]
    defines a left $T$-module structure on $T\otimes_R M$.

    This proves \textup{(2)}, and hence also \textup{(1)}.
\end{proof}

\begin{proposition}
    Let $I$ be an ideal of $R$. Let $M$ be an $R$-module. Then the morphism of $R$-modules:
    \[
        M \to (R/I) \otimes_R M
    \]
    defined by
    \[
        m \mapsto 1_{R/I} \otimes m
    \]
    induces an isomorphism of $R$-modules:
    \[
        M/(IM) \cong (R/I) \otimes_R M
    \]
    as $R/I$-modules.
\end{proposition}

\begin{proof}
    Define
    \[
        \varphi:M\longrightarrow (R/I)\otimes_R M,\qquad m\longmapsto 1_{R/I}\otimes m.
    \]
    This is an $R$-module homomorphism. For any $a\in I$ and $m\in M$, we have
    \[
        \varphi(am)=1_{R/I}\otimes am=(1_{R/I}\cdot a)\otimes m=(a+I)\otimes m=0,
    \]
    since $a+I=0$ in $R/I$. Hence $\varphi(IM)=0$, so $IM\subseteq \ker(\varphi)$.

    Therefore, by the universal property of the quotient,$\varphi$ induces an $R$-module
    homomorphism
    \[
        \overline{\varphi}:M/IM\longrightarrow (R/I)\otimes_R M,
        \qquad m+IM\longmapsto 1_{R/I}\otimes m.
    \]

    Now define
    \[
        \beta:(R/I)\times M\longrightarrow M/IM,\qquad (r+I,m)\longmapsto rm+IM.
    \]
    We first check that $\beta$ is well defined. If $r+I=r'+I$, then $r-r'\in I$, so
    \[
        (r-r')m\in IM,
    \]
    hence
    \[
        rm+IM=r'm+IM.
    \]
    Also,$\beta$ is $R$-balanced, since for any $a\in R$,
    \[
        \beta((r+I)a,m)=\beta(ra+I,m)=ram+IM=\beta(r+I,am).
    \]
    Thus, by the universal property of the tensor product, there exists a unique
    homomorphism
    \[
        \psi:(R/I)\otimes_R M\longrightarrow M/IM
    \]
    such that
    \[
        \psi((r+I)\otimes m)=rm+IM
        \qquad\text{for all }r\in R,\ m\in M.
    \]

    We now show that $\overline{\varphi}$ and $\psi$ are inverse to each other.

    For any $m\in M$,
    \[
        (\psi\circ \overline{\varphi})(m+IM)
        =\psi(1_{R/I}\otimes m)
        =1m+IM
        =m+IM.
    \]
    So
    \[
        \psi\circ \overline{\varphi}=\mathrm{id}_{M/IM}.
    \]

    For any elementary tensor $(r+I)\otimes m\in (R/I)\otimes_R M$,
    \[
        (\overline{\varphi}\circ \psi)\bigl((r+I)\otimes m\bigr)
        =\overline{\varphi}(rm+IM)
        =1_{R/I}\otimes rm
        =(r+I)\otimes m.
    \]
    Since elementary tensors generate $(R/I)\otimes_R M$, it follows that
    \[
        \overline{\varphi}\circ \psi=\mathrm{id}_{(R/I)\otimes_R M}.
    \]

    Hence $\overline{\varphi}$ is an isomorphism:
    \[
        M/IM\cong (R/I)\otimes_R M.
    \]

    Finally, both modules carry natural $R/I$-module structures, and the maps $\overline{\varphi}$ and $\psi$ are $R/I$-linear. Therefore this is an isomorphism
    of $R/I$-modules.
\end{proof}

